\documentclass[11pt]{article}
\usepackage{icmlko}
\begin{document}

\title{정렬 수집이 늘 절단은 아니다 --- 바닥이 어디 떠 있나로 갈린다}
\author{월드모델 실험실}
\date{노트 375 · 2026년 8월 1일}
\maketitle

\begin{abstract}
노트 374가 ``수집 정렬이 과제를 정한 \textbf{네 번째} 도메인''이라 적으며
도서 · 모바일 · 만화 · 펀딩을 한 묶음으로 놨다. \textbf{묶으면 안 됐다.}
라벨의 \emph{바닥}을 재면 넷이 갈린다 --- 최소값을 1분위로 나누면
세계애니 $0.98$ · 만화 $0.94$ 로 바닥에 몰려 있고, 웹툰 $0.04$ · 도서
$0.17$ · 펀딩 $0.47$ · 게임 $0.61$ 은 아래로 열려 있다(모바일은 $1.00$
이지만 최소가 \textbf{1} 이라 자연 바닥이지 절단이 아니다).
\textbf{정렬해서 모아도 충분히 깊이 내려가면 절단이 안 생긴다.}

그리고 지도를 그리다 못 보던 것이 나왔다. \textbf{제일 심하게 잘린 것은
세계애니다} --- 최소 popularity가 \textbf{19,533} 이고 중앙이 49,618 이라
\emph{도메인 전체가 바닥에서 중앙까지 2.5배}밖에 안 된다. 가중 10\% ·
$\rho=0.5567$ 로 판에서 셋째로 높은 도메인인데, 그 수는 ``애니 중에서''가
아니라 \textbf{``인기 상위 3,000편 중에서''} 매긴 순위다. 스무 편 넘게
이 도메인을 인용했지만 아무도 그 단서를 안 달았다.
\end{abstract}

\section{바닥을 잰다}

정렬해서 모으면 위에서부터 채워지므로 \emph{바닥}에 자국이 남는다.
최소값이 1분위와 붙어 있으면(비 $\to 1$) 잘린 것이고, 멀면 아래로 열려
있다. 다만 최소가 \emph{자연 바닥}(리뷰 1개, 후원자 1명)에 닿아 있으면
잘린 게 아니다.

\begin{center}\includegraphics[width=\linewidth]{figs/floor.pdf}\end{center}

\begin{center}
\begin{tabular}{lrrrrrl}
\hline
도메인 & $n$ & 최소 & 1분위 & 중앙 & 최소$\div$1분위 & 판정 \\
\hline
\textbf{세계애니} & 2,988 & \textbf{19,533} & 19,833 & 49,618 & \textbf{0.98} & \textbf{절단} \\
\textbf{만화} & 2,983 & \textbf{632} & 672 & 7,864 & \textbf{0.94} & \textbf{절단} \\
모바일 & 2,000 & 1 & 1 & 352 & 1.00 & 자연 바닥 \\
게임 & 482 & 45 & 73 & 24,292 & 0.61 & 열림 \\
펀딩 & 400 & 9 & 19 & 139 & 0.47 & 열림 \\
도서 & 396 & 10 & 60 & 8,860 & 0.17 & 열림 \\
웹툰 & 2,817 & 19 & 474 & 55,208 & 0.04 & 열림 \\
\hline
\end{tabular}
\end{center}

\textbf{노트 374의 묶음이 깨진다.} 도서는 베스트셀러 목록에서 모았고
모바일은 인기 차트에서 모았는데 \emph{둘 다 안 잘렸다} --- 도서는 분야마다
3,000위까지 훑고 모바일은 새 앱이 차트에 올라 리뷰 1개짜리까지 들어온다.
\textbf{정렬은 조건이고, 잘리느냐는 얼마나 깊이 내려갔느냐가 정한다.}

\section{수집 편향이 세 가지다}

\begin{center}
\begin{tabular}{lll}
\hline
종류 & 자국 & 예 \\
\hline
\textbf{절단} & 바닥이 자연 바닥 위에 떠 있다 & 세계애니 · 만화 \\
\textbf{구역 편향} & 바닥은 열렸는데 표본이 한쪽에서 왔다 & 펀딩(위 19\%) \\
\textbf{열림} & 바닥이 자연 바닥까지 내려간다 & 도서 · 모바일 · 게임 · 웹툰 \\
\hline
\end{tabular}
\end{center}

\textbf{셋을 가르는 검사가 다르다.} 절단은 \emph{바닥}으로 보이고, 구역
편향은 바닥으로 안 보인다 --- 노트 374의 펀딩은 최소 9명까지 내려가는데도
쪽 50\textasciitilde{}753 에서만 왔다. 그건 \emph{수집 순서와 라벨의 상관}
으로만 보였다($\rho=-0.3666$). \textbf{두 검사를 다 돌려야 한다.}

\section{세계애니가 제일 심하다}

\begin{center}
\begin{tabular}{lr}
\hline
세계애니 popularity 분위 & \\
\hline
최소 & \textbf{19,533} \\
1\% & 19,832 \\
25\% & 30,449 \\
50\% & 49,618 \\
75\% & 98,169 \\
최대 & 1,035,081 \\
\hline
\end{tabular}
\end{center}

\textbf{바닥에서 중앙까지 2.5배, 바닥에서 1분위까지 1.5\%.} 위쪽은 최대가
중앙의 21배로 길게 열려 있는데 아래가 칼로 자른 듯 끊긴다 --- 인기순
상위 $\sim$3,000편을 받아 온 자국이다.

\textbf{그런데 그 수집 스크립트가 저장소에 없다.} \texttt{wanime\_records.json}
을 쓰는 코드는 아홉 군데가 읽기만 하고 만드는 쪽이 없다 --- 노트 365의
\texttt{media\_push\_fill} 에 이어 \textbf{두 번째}다. 그래서 정렬은
\emph{바닥 모양으로 추론}할 수밖에 없었다.

\textbf{읽는 법이 바뀐다.} 세계애니 $\rho=0.5567$(95\% $[0.462,\,0.646]$,
가중 10\%)은 ``세계 애니 중에서''가 아니라 \textbf{``AniList 인기 상위
3,000편 중에서''} 매긴 순위다. 노트 284가 시장팝업에 붙인 단서
(``결과로 선택된 모집단이다'')와 같은 종류이고, 노트 370이 만화에 붙인
것과도 같다. \textbf{세계애니에는 아무도 안 붙였다} --- 스무 편 넘게
인용하면서.

\section{규약에 더하는 줄}

\begin{itemize}
\item \textbf{정렬 수집을 봤으면 바닥을 잰다.} 최소 $\div$ 1분위가 1 에
가깝고 최소가 자연 바닥 위에 있으면 절단이다.
\item \textbf{절단과 구역 편향은 다른 검사로 잡힌다} --- 바닥, 그리고
수집 순서 $\leftrightarrow$ 라벨 상관. 둘 다 돌린다.
\item \textbf{모집단 단서는 도메인 문서에 박는다.} 시장팝업(284) ·
만화(370)에는 있고 세계애니에는 없었다.
\end{itemize}

\section{남은 것}

\begin{center}
\begin{tabular}{ll}
\hline
갈래 & 지금 아는 것 \\
\hline
\textbf{균일 펀딩 2,500건} & 돌고 있다(노트 374). 다 모이면 짝으로 잰다 \\
세계애니를 아래로 넓히나 & 바닥 19,533 아래를 받으면 도메인이 넓어진다. \\
 & 다만 만화 KR 처럼 부호가 뒤집힐 수 있다(노트 372) --- \\
 & 넣기 전에 부분집합으로 따로 채점해 본다 \\
팝업 실측 41행 & 노트 360 · 364 · 366 \\
\hline
\end{tabular}
\end{center}

둘째 줄이 이 노트가 여는 것이다. 세계애니는 가중 10\% 에 유보 300행이라
넓히면 판이 크게 움직인다 --- 그리고 노트 372가 그 위험을 이미 이름
붙여 놨다.

\end{document}
